\sloppy
\emergencystretch=5em
\hbadness=10000
\vbadness=10000
\tolerance=2000

\providecolor{codebg}{RGB}{248,248,248}
\providecolor{codeframe}{RGB}{200,200,200}
\providecolor{keyword}{RGB}{0,0,180}
\providecolor{string}{RGB}{163,21,21}
\providecolor{comment}{RGB}{0,128,0}

\lstset{
	backgroundcolor=\color{codebg},
	frame=single,
	rulecolor=\color{codeframe},
	language=Python,
	basicstyle=\ttfamily\scriptsize,
	keywordstyle=\color{keyword}\bfseries,
	stringstyle=\color{string},
	commentstyle=\color{comment}\itshape,
	showstringspaces=false,
	breaklines=true,
	breakatwhitespace=true,
	numbers=left,
	numberstyle=\tiny\color{gray},
	numbersep=8pt,
	tabsize=4,
	literate=
	{ف}{{{\bfseries\itshape ف}}}1
	{ی}{{{\bfseries\itshape ی}}}1
}

\section{نمای کلی و هدف ماژول}

ماژول \lr{\texttt{lp\_validation.py}} یک اسکریپت اعتبارسنجی \emph{بهینه‌سازی پارامترهای پروتکل} است که نقش یک \lr{validator} مستقل را در فریم‌ورک شبیه‌سازی \lr{QKD} ایفا می‌کند. این ماژول به‌طور خاص برای پاسخ به پرسش بنیادین زیر طراحی شده است: \emph{آیا حلّال بهینه‌سازی پارامترهای پروتکل \lr{Lim-2014} (\lr{LP Solver}) که در ماژول \lr{\texttt{qkd.proofs.optimization}} پیاده‌سازی شده، در یافتن پارامترهای بهینه (\lr{optimal}) به‌طور قابل‌اتکا عمل می‌کند و آیا خروجی آن حداقل به‌اندازه پارامترهای سخت‌کدشده (\lr{hardcoded}) و مرجع، نرخ کلید امن (\lr{Secure Key Rate} یا \lr{SKR}) تولید می‌نماید؟} این پرسش از آنجا اهمیت می‌یابد که در پروتکل‌های \lr{decoy-state BB84}، انتخاب پارامترهای شدت پالس (\lr{$\mu$} و \lr{$\nu$}) و احتمال‌های تخصیص (\lr{$p_s$} و \lr{$p_d$}) تأثیر مستقیمی بر نرخ کلید نهایی دارد و یک بهینه‌ساز ناقص می‌تواند منجر به انتخاب پارامترهای زیربهینه (\lr{suboptimal}) شود که نرخ کلیدی کم‌تر از حد قابل دستیابی تولید می‌کنند.

هدف اصلی این ماژول، \emph{اعتبارسنجی تطبیقی} (\lr{comparative validation}) حلّال بهینه‌سازی در برابر یک مرجع سخت‌کدشده است. به‌عبارت دیگر، این ماژول یک نقطه عملیاتی مرجع با پارامترهای ثابت \lr{$\mu = 0.5$}، \lr{$\nu = 0.1$}، \lr{$q_x = 0.5$}، \lr{$p_s = 0.5$} و \lr{$p_d = 0.25$} (که در ادبیات \lr{decoy-state QKD} به‌عنوان مقادیر کلاسیک و شناخته‌شده هستند) تعریف می‌کند، سپس حلّال بهینه‌سازی را فراخوانی می‌نماید و در صورتی که \lr{SKR} بهینه حداقل برابر با \lr{SKR} مرجع باشد، آزمون موفق تلقی می‌گردد. این رویکرد، الگوی \lr{necessary condition testing} است: موفقیت آزمون تضمین می‌کند که حلّال حداقل به اندازه مرجع عمل می‌کند؛ هرچند ممکن است پارامترهای بهتر نیز یافت کند که در آن صورت، آزمون با حاشیه مثبت موفق می‌شود.

ساختار کلی این ماژول شامل سه تابع است. تابع \lr{\texttt{H2}} آنتروپی باینری شانون را با برش عددی (\lr{numerical clipping}) پیاده‌سازی می‌کند تا از خطای \lr{math domain error} در نقاط \lr{$p = 0$} و \lr{$p = 1$} جلوگیری نماید. تابع \lr{\texttt{analytical\_skr}} تابع هدف تحلیلی را محاسبه می‌کند که از حلّال بهینه‌سازی به‌عنوان \lr{cost function} استفاده می‌گردد؛ این تابع، پنج پارامتر پروتکل را به‌عنوان ورودی می‌گیرد و \lr{SKR} مجانبی را بر اساس فرمول \lr{GLLP} (\lr{Gottesman-Lo-Lütkenhaus-Preskill}) بازمی‌گرداند. تابع \lr{\texttt{run\_lp\_test}} آزمون اصلی است که در یک فاصله مشخص، حلّال را فراخوانی کرده و نتایج را مقایسه می‌نماید.

ویژگی برجسته این ماژول، \emph{استفاده از تابع callback} برای تزریق تابع هدف است. به‌جای آنکه حلّال بهینه‌سازی به‌صورت سخت‌کدشده به فرمول \lr{SKR} وابسته باشد، تابع \lr{\texttt{cost\_callback}} به‌عنوان پارامتر به حلّال ارسال می‌شود. این الگو، الگوی \lr{dependency injection} است و امکان جایگزینی تابع هدف با توابع دیگر (مثلاً \lr{SKR} متناهی-\lr{key} یا توابع هزینه پیچیده‌تر) را بدون تغییر در حلّال فراهم می‌سازد. این طراحی، جداسازی واضحی بین منطق بهینه‌سازی و تابع هدف برقرار می‌کند و ماژولار بودن سیستم را افزایش می‌دهد.

در سطح عملیاتی، این ماژول سه آزمون در فواصل مختلف ($0$، $50$ و $100$ کیلومتر) اجرا می‌کند تا عملکرد حلّال در رژیم‌های مختلف کانال (بدون کاهیدگی، کاهیدگی متوسط و کاهیدگی شدید) را راستی‌آزمایی نماید. در فاصله $L = 0$ کیلومتر، ضریب انتقال کانال برابر با $1$ است و حلّال باید پارامترهایی با \lr{SKR} نزدیک به حد نظری پیدا کند؛ در فاصله $L = 50$ کیلومتر، ضریب انتقال به $\eta_{\text{channel}} \approx 0.1$ کاهش می‌یابد و حلّال باید با کاهش $\mu$ برای جبران کاهیدگی، نرخ کلید را بهینه نماید؛ در فاصله $L = 100$ کیلومتر، با ضریب انتقال $\eta_{\text{channel}} \approx 0.01$، حلّال در رژیم نزدیک به آستانه بحرانی \lr{BB84} (\lr{$Q \approx 0.11$}) عمل می‌کند و انتخاب پارامترها به‌طور حساس به نتیجه تأثیر می‌گذارد. این سه آزمون، سه رژیم فیزیکی متفاوت را پوشش می‌دهند و تضمین می‌کنند که حلّال در طیف وسیعی از شرایط عملیاتی به‌طور قابل‌اتکا عمل می‌کند.

نکته مهم دیگر آن است که این ماژول \emph{یک تابع هدف تحلیلی مستقل} برای اعتبارسنجی استفاده می‌کند، نه همان تابعی که حلّال در داخل خود استفاده می‌نماید. این جداسازی، الگوی \lr{independent oracle} است و از خطای \lr{circular validation} جلوگیری می‌کند؛ یعنی اگر حلّال از یک تابع هدف باگ‌دار استفاده کند، این باگ در اعتبارسنجی نیز منتقل نشود و شناسایی گردد. تابع \lr{\texttt{analytical\_skr}} در این ماژول، مستقل از حلّال پیاده‌سازی شده و به‌عنوان یک \lr{oracle} عمل می‌کند که صحت خروجی حلّال را راستی‌آزمایی می‌نماید.

\section{مبانی تئوری و فیزیکی}

\subsection{پروتکل Decoy-State و نقش پارامترهای شدت}

در پروتکل \lr{decoy-state BB84}، فرستنده (\lr{Alice}) به‌جای ارسال پالس با شدت ثابت، به‌طور تصادفی پالس‌ها را در سه دسته با شدت‌های مختلف توزیع می‌کند: سیگنال با شدت $\mu$، دکوی ضعیف با شدت $\nu$ (\lr{$\nu < \mu$}) و واکوئم با شدت صفر. هدف از این توزیع، آن است که حمله \lr{Photon-Number Splitting} (\lr{PNS}) که علیه منبع \lr{Weak Coherent Pulse} (\lr{WCP}) با توزیع \lr{Poisson} قابل اعمال است، خنثی شود. در این حمله، حاملهور (\lr{eavesdropper}) با تقسیم پالس‌های چندفوتونی، یک کپی از فوتون را نگه می‌دارد و بقیه را به گیرنده ارسال می‌کند؛ اما در حضور دکوی، حاملهور نمی‌تواند تشخیص دهد که کدام پالس سیگنال و کدام دکوی است و بنابراین، حمله \lr{PNS} قابل اعمال نیست.

پارامترهای کلیدی این پروتکل عبارت‌اند از: $q_x$ (\lr{sifting factor}، کسر پالس‌هایی که پایه‌های \lr{Alice} و \lr{Bob} هم‌خوان هستند)، $p_s$ (احتمال تخصیص پالس به سیگنال)، $p_d$ (احتمال تخصیص پالس به دکوی)، $p_v = 1 - p_s - p_d$ (احتمال تخصیص به واکوئم)، $\mu$ (شدت میانگین سیگنال) و $\nu$ (شدت میانگین دکوی). انتخاب این پارامترها، یک مسئله بهینه‌سازی غیرخطی است که هدف آن، حداکثرسازی نرخ کلید امن است. این بهینه‌سازی، به‌دلیل وابستگی پیچیده \lr{SKR} به این پارامترها (از طریق استخراج $Y_1$ و محاسبه $e_1$ و $E_\mu$)، به‌سادگی قابل حل نیست و نیاز به حلّال‌های عددی نظیر \lr{SLSQP} یا \lr{Nelder-Mead} دارد.

\subsection{آنتروپی باینری شانون و نقش آن در امنیت}

آنتروپی باینری شانون که با $H_2(p)$ نشان داده می‌شود، یکی از مفاهیم بنیادین نظریه اطلاعات است که در رمزنگاری کوانتومی نقشی مرکزی ایفا می‌کند. این کمیت، میزان عدم‌قطعیت متغیر تصادفی باینری با توزیع \lr{Bernoulli} را اندازه می‌گیرد:

\begin{equation}
	H_2(p) = -p \log_2(p) - (1 - p) \log_2(1 - p)
\end{equation}

در این رابطه، $p \in [0, 1]$ احتمال وقوع یکی از دو حالت است. تابع $H_2(p)$ در $p = 0.5$ به مقدار بیشینه $1$ می‌رسد و در نقاط $p = 0$ و $p = 1$ به صفر میل می‌کند. در رمزنگاری کوانتومی، $H_2(p)$ در دو نقطه بحرانی ظاهر می‌شود: نخست در محاسبه \lr{privacy amplification} که حجم اطلاعات نشت‌کرده به حاملهور را تعیین می‌کند (\lr{$H_2(E_\mu)$} و \lr{$H_2(e_1)$})، و دوم در محاسبه ضریب \lr{privacy amplification} در فرمول \lr{GLLP}.

پیاده‌سازی تابع $H_2$ در کد، با برش عددی در نقاط $p = 0$ و $p = 1$ همراه است، زیرا در این نقاط، عبارت $\log_2(p)$ به $-\infty$ میل می‌کند. در محاسبات عددی، این حالت می‌تواند منجر به \lr{NaN} یا \lr{ValueError} شود؛ بنابراین پیاده‌سازی پیش از فراخوانی \lr{\texttt{math.log2}}، شرایط $p \leq 0$ یا $p \geq 1$ را بررسی کرده و در این موارد، مقدار $0$ را برمی‌گرداند. این برش با حد گسسته \lr{L'Hôpital} هم‌خوان است:

\begin{equation}
	\lim_{p \to 0^+} p \log_2 p = \lim_{p \to 0^+} \frac{\log_2 p}{1/p} = \lim_{p \to 0^+} \frac{-p}{\ln 2} = 0
\end{equation}

\subsection{کاهیدگی فیبر و قانون Beer-Lambert}

انتقال فوتون در فیبر نوری تک‌حالت در طول‌موج مخابراتی $1550$ نانومتر با قانون \lr{Beer-Lambert} مدل می‌شود. برای یک فیبر به طول $L$ کیلومتر با ضریب کاهیدگی خطی $\alpha_{\text{dB/km}}$، کاهیدگی کل بر حسب دسی‌بل برابر با $\alpha_{\text{dB/km}} \cdot L$ است و تبدیل آن به ضریب انتقال خطی از رابطه زیر پیروی می‌کند:

\begin{equation}
	\eta_{\text{channel}}(L, \alpha_{\text{dB/km}}) = 10^{-\frac{\alpha_{\text{dB/km}} \cdot L}{10}}
\end{equation}

در این پیاده‌سازی، $\alpha_{\text{dB/km}} = 0.2$ (مقدار استاندارد فیبر تک‌حالت در $1550$ نانومتر) به‌کار رفته است. در فاصله $L = 50$ کیلومتر، $\eta_{\text{channel}} = 10^{-1} = 0.1$ و در فاصله $L = 100$ کیلومتر، $\eta_{\text{channel}} = 10^{-2} = 0.01$ می‌شود. ضریب انتقال کلی سیستم شامل بازده آشکارساز $\eta_{\text{det}}$ نیز می‌شود:

\begin{equation}
	\eta_{\text{total}}(L) = \eta_{\text{det}} \cdot \eta_{\text{channel}}(L, \alpha_{\text{dB/km}})
\end{equation}

در این پیاده‌سازی، $\eta_{\text{det}} = 0.15$ که متناظر با آشکارساز \lr{InGaAs SPAD} در دمای $-50$ درجه سانتی‌گراد است. بنابراین در $L = 50$ کیلومتر، $\eta_{\text{total}} = 0.015$ می‌شود.

\subsection{Gain سیگنال و QBER سیگنال}

برای منبع \lr{WCP} با توزیع \lr{Poisson} با شدت میانگین $\mu$، \lr{Gain} سیگنال (یعنی احتمال آنکه یک پالس منجر به ثبت کلیک در آشکارساز شود) برابر است با:

\begin{equation}
	Q_\mu(\mu, \eta, Y_0) = 1 - (1 - Y_0) \cdot e^{-\eta \mu}
\end{equation}

در این رابطه، $Y_0$ نرخ \lr{dark count} به ازای هر پالس، $\eta$ ضریب انتقال کلی سیستم و $\mu$ شدت میانگین سیگنال است. عبارت $e^{-\eta \mu}$ احتمال آنکه هیچ فوتونی شناسایی نشود و عامل $(1 - Y_0)$ احتمال آنکه دارک کانت نیز رخ ندهد را نشان می‌دهد. \lr{QBER} سیگنال، یعنی کسر کلیک‌های اشتباه در میان کلیک‌های ثبت‌شده، از رابطه زیر به دست می‌آید:

\begin{equation}
	E_\mu = \frac{0.5 \cdot Y_0 + e_d \cdot (1 - e^{-\eta \mu})}{Q_\mu}
\end{equation}

در این رابطه، $e_d$ خطای ذاتی سیستم (\lr{misalignment error}) است. صورت این کسر از دو بخش تشکیل شده است: $0.5 \cdot Y_0$ که نشان‌دهنده کلیک‌های اشتباه ناشی از دارک کانت است (چون دارک کانت کاملاً تصادفی است، نصف آن در پایه درست و نصف در پایه اشتباه رخ می‌دهد) و $e_d \cdot (1 - e^{-\eta \mu})$ که نشان‌دهنده کلیک‌های اشتباه ناشی از فوتون‌های واقعی با احتمال خطای ذاتی $e_d$ است.

\subsection{Gain دکوی و استخراج Y1 به روش Lo2005}

\lr{Gain} دکوی ضعیف به‌صورت مشابه محاسبه می‌شود:

\begin{equation}
	Q_\nu(\nu, \eta, Y_0) = 1 - (1 - Y_0) \cdot e^{-\eta \nu}
\end{equation}

مرحله کلیدی در \lr{decoy-state QKD}، استخراج کران پایین برای $Y_1$ (احتمال آنکه یک پالس تک‌فوتونی منجر به کلیک شود) از مشاهدات $Q_\mu$ و $Q_\nu$ است. این کران، که در کار مرجع \lr{Lo-Ma-Chen (2005)} به‌تفصیل استخراج شده، در این ماژول به فرم زیر پیاده‌سازی شده است:

\begin{equation}
	Y_1^{L} = \frac{\mu}{\mu \nu - \nu^2} \left( Q_\nu \cdot e^{\nu} - \frac{\nu^2}{\mu^2} \cdot Q_\mu \cdot e^{\mu} \right)
\end{equation}

این فرمول، یک نسخه بازآرایی‌شده از فرمول کلاسیک \lr{Lo2005} است. ضریب $\frac{\mu}{\mu \nu - \nu^2} = \frac{\mu}{\nu(\mu - \nu)}$ در صورتی که $\nu \to \mu$ باشد، به بی‌نهایت میل می‌کند؛ بنابراین در پیاده‌سازی، شرط $\left| \mu - \nu \right| < 10^{-9}$ بررسی می‌شود و در این حالت، تابع مقدار صفر را برمی‌گرداند. در این پیاده‌سازی، یک اصلاح فیزیکی مهم نیز اعمال شده است: مقدار نهایی $Y_1$ با \lr{\texttt{max(Y1\_L, Y0)}} محاسبه می‌شود تا تضمین گردد که $Y_1$ هرگز کم‌تر از نرخ دارک کانت $Y_0$ نشود. این برش، بازتاب‌دهنده این واقعیت فیزیکی است که حتی در غیاب فوتون واقعی، دارک کانت می‌تواند منجر به کلیک شود و بنابراین، بازده تک‌فوتونی نمی‌تواند کم‌تر از $Y_0$ باشد.

\subsection{نرخ خطای تک‌فوتونی و بازده تک‌فوتونی}

نرخ خطای تک‌فوتونی $e_1$، یعنی کسر کلیک‌های اشتباه در میان پالس‌های تک‌فوتونی، از رابطه زیر به دست می‌آید:

\begin{equation}
	e_1 = \frac{0.5 \cdot Y_0 + e_d \cdot \eta}{Y_1}
\end{equation}

این رابطه مشابه $E_\mu$ است، با این تفاوت که به‌جای احتمال کلیک فوتون واقعی $1 - e^{-\eta \mu}$، احتمال کلیک تک‌فوتونی $\eta$ قرار دارد. در پیاده‌سازی، شرط $Y_1 > 0$ بررسی می‌شود تا از تقسیم بر صفر جلوگیری شود؛ در صورت $Y_1 = 0$، مقدار $e_1 = 0.5$ گرفته می‌شود که نشان‌دهنده حالت کاملاً تصادفی است (یعنی در غیاب سیگنال واقعی، نتیجه اندازهگیری کاملاً تصادفی و با احتمال برابر برای هر بیت است).

\lr{Gain} تک‌فوتونی، یعنی کسری از کلیک‌های کل که از پالس‌های تک‌فوتونی ناشی می‌شوند، از رابطه زیر به دست می‌آید:

\begin{equation}
	Q_1 = \mu \cdot e^{-\mu} \cdot Y_1
\end{equation}

این رابطه از ضرب احتمال آنکه پالس تک‌فوتونی باشد (\lr{$P(1 | \mu) = \mu e^{-\mu}$}) در بازده تک‌فوتونی $Y_1$ به دست می‌آید.

\subsection{نرخ کلید امن مجانبی (GLLP)}

نرخ کلید امن مجانبی (\lr{asymptotic secure key rate}) از فرمول کلاسیک \lr{GLLP} به دست می‌آید:

\begin{equation}
	R = q_x \cdot \left[ Q_1 \cdot (1 - H_2(e_1)) - Q_\mu \cdot f_{\text{EC}} \cdot H_2(E_\mu) \right]
\end{equation}

این فرمول از دو بخش تشکیل شده است. بخش اول $Q_1 \cdot (1 - H_2(e_1))$ نشان‌دهنده حجم اطلاعات مشترک بین فرستنده و گیرنده از پالس‌های تک‌فوتونی پس از \lr{privacy amplification} است؛ عبارت $1 - H_2(e_1)$ ضریب \lr{privacy amplification} است که حجم اطلاعات نشت‌کرده به حاملهور را کم می‌کند. بخش دوم $Q_\mu \cdot f_{\text{EC}} \cdot H_2(E_\mu)$ نشان‌دهنده حجم اطلاعات فاش‌شده در مرحله \lr{error correction} است؛ $f_{\text{EC}} \geq 1$ بازدهی \lr{error correction} است که در عمل $f_{\text{EC}} \approx 1.16$ برای کد \lr{LDPC} بهینه است. در این پیاده‌سازی، $f_{\text{EC}} = 1.16$ به‌کار رفته است که با مقادیر عملی در پیاده‌سازی‌های \lr{QKD} هم‌خوان است.

تفاوت این فرمول با حالت \lr{finite-key} (مثل اثبات \lr{Lim-2014}) آن است که در اینجا از حد مجانبی $N \to \infty$ استفاده می‌شود و جریمه‌های آماری نظیر \lr{Hoeffding bound} و \lr{composable security} اعمال نمی‌گردند. به همین دلیل، $R$ در اینجا همیشه بزرگ‌تر یا مساوی با نرخ کلید متناهی است و در عمل به‌عنوان \emph{کران بالای} نرخ کلید قابل دستیابی عمل می‌کند. با این حال، تطابق حلّال بهینه‌سازی با این کران بالا، شرط لازم — هرچند کافی نیست — برای صحت پیاده‌سازی است.

\subsection{قیدهای فیزیکی و ریاضی}

در تابع \lr{\texttt{analytical\_skr}}، چندین قید فیزیکی و ریاضی اعمال شده است که در صورت نقض، تابع مقدار صفر را برمی‌گرداند. این قیدها عبارت‌اند از:

نخست، قید \lr{$p_s + p_d \geq 1$}: این قید تضمین می‌کند که مجموع احتمال‌های تخصیص سیگنال و دکوی از $1$ تجاوز نکند؛ در غیر این صورت، احتمال واکوئم $p_v = 1 - p_s - p_d$ منفی می‌شود که از لحاظ فیزیکی بی‌معناست. در صورت نقض، تابع مقدار صفر برمی‌گرداند.

دوم، قید \lr{$\nu \geq 0.9 \cdot \mu$}: این قید تضمین می‌کند که شدت دکوی به‌طور قابل‌توجهی کم‌تر از شدت سیگنال باشد؛ در غیر این صورت، استخراج $Y_1$ از روش \lr{Lo2005} به‌دلیل نزدیک بودن $\mu$ و $\nu$ به تکینگی (\lr{singularity}) می‌رسد و نتایج غیرقابل اعتماد تولید می‌کند. در عمل، $\nu$ معمولاً در بازه $0.1 \mu$ تا $0.5 \mu$ انتخاب می‌شود.

سوم، قید \lr{$\left| \mu - \nu \right| < 10^{-9}$}: این قید حالت تکین را شناسایی می‌کند و در صورت وقوع، تابع مقدار صفر برمی‌گرداند.

این قیدها، الگوی \lr{defensive programming} را به نمایش می‌گذارند: ورودی‌ها پیش از محاسبه اعتبارسنجی می‌شوند و در صورت غیرمعتبر بودن، مقدار پیش‌فرض امن (صفر) بازگردانده می‌شود. این الگو در توابع بهینه‌سازی حیاتی است، چون حلّال ممکن است در جریان جستجوی فضای پارامتر، به مقادیر غیرفیزیکی برسد و بدون این قیدها، تابع هدف می‌تواند مقادیر غیرمتعارف یا \lr{NaN} تولید کند.

\section{تحلیل معماری و ساختار داده‌ها}

\subsection{ساختار کلی ماژول و الگوی Module-Level Functions}

این ماژول از الگوی \lr{module-level functions} پیروی می‌کند، یعنی به‌جای تعریف کلاس‌ها، از توابع مستقل در سطح ماژول استفاده می‌نماید. این انتخاب طراحی، بازتاب‌دهنده ماهیت \lr{procedural} این ماژول است: توابع تحلیلی (\lr{\texttt{H2}} و \lr{\texttt{analytical\_skr}}) توابع ریاضی خالص (\lr{pure functions}) هستند که ورودی‌هایشان را به خروجی‌هایشان نگاشت می‌کنند و هیچ \lr{state} داخلی ندارند. تابع \lr{\texttt{run\_lp\_test}} تنها تابع دارای \lr{side-effect} است که خروجی چاپ می‌کند و حلّال را فراخوانی می‌نماید. این جداسازی بین توابع خالص و توابع دارای \lr{side-effect}، یکی از اصول مهندسی نرم‌افزار بالغ است و در الگوهای طراحی نظیر \lr{functional core, imperative shell} تبیین می‌شود.

\subsection{تنظیم مسیر و Path Bootstrap}

در ابتدای ماژول، الگوی \lr{path bootstrap} پیاده‌سازی شده است:

\begin{LTR}
	\begin{lstlisting}[language=Python]
		sys.path.insert(0, os.path.abspath(os.path.dirname(__file__)))
	\end{lstlisting}
\end{LTR}

این خط، مسیر دایرکتوری حاوی اسکریپت را به ابتدای \lr{\texttt{sys.path}} اضافه می‌کند. این کار تضمین می‌نماید که ماژول بدون توجه به محل اجرای آن (از \lr{project root} یا از مسیر دیگر)، بتواند ماژول‌های \lr{\texttt{qkd.proofs.optimization}} و \lr{\texttt{qkd.params}} را وارد نماید. این الگو در اسکریپت‌هایی که به‌صورت مستقل اجرا می‌شوند و نیاز به واردسازی ماژول‌های پروژه دارند، استاندارد است.

\subsection{واردسازی توابع از فریم‌ورک}

دو واردسازی کلیدی در این ماژول انجام می‌شود:

\begin{LTR}
	\begin{lstlisting}[language=Python]
		from qkd.proofs.optimization import optimize_lim2014_finite_key
		from qkd.params import load_lim2014_dedicated_params
	\end{lstlisting}
\end{LTR}

نخست، تابع \lr{\texttt{optimize\_lim2014\_finite\_key}} از ماژول \lr{\texttt{qkd.proofs.optimization}} وارد می‌شود. این تابع، حلّال بهینه‌سازی پارامترهای پروتکل \lr{Lim-2014} است که از الگوریتم \lr{Nelder-Mead} یا \lr{SLSQP} برای یافتن پارامترهای بهینه (\lr{$q_x$}، \lr{$p_s$}، \lr{$p_d$}، \lr{$\mu$}، \lr{$\nu$}) استفاده می‌کند. دوم، تابع \lr{\texttt{load\_lim2014\_dedicated\_params}} از ماژول \lr{\texttt{qkd.params}} وارد می‌شود. این تابع، پارامترهای اختصاصی پروتکل \lr{Lim-2014} (مانند \lr{$N$}، \lr{$\varepsilon_{\text{PE}}$}، \lr{$\varepsilon_{\text{EC}}$} و سایر پارامترهای امنیتی) را برای یک فاصله و تعداد بیت مشخص بارگذاری می‌نماید.

این وابستگی‌ها نشان‌دهنده آن است که ماژول اعتبارسنجی، مستقیماً به موتور بهینه‌سازی فریم‌ورک متصل است و همان مسیر اجرای واقعی را آزمون می‌کند. این الگو، تضمین می‌کند که آزمون، رفتار واقعی موتور را در شرایط عملیاتی منعکس می‌کند و نه رفتار آزمایشگاهی مصنوعی.

\subsection{تحلیل تابع H2}

تابع \lr{\texttt{H2(p)}} یک تابع ریاضی خالص است که آنتروپی باینری شانون را محاسبه می‌نماید:

\begin{LTR}
	\begin{lstlisting}[language=Python]
		def H2(p):
		if p <= 0 or p >= 1: return 0.0
		return -p * math.log2(p) - (1 - p) * math.log2(1 - p)
	\end{lstlisting}
\end{LTR}

در خط اول، شرط $p \leq 0$ یا $p \geq 1$ بررسی می‌شود. این شرط شامل نقاط انتهایی $p = 0$ و $p = 1$ نیز می‌شود که در آن‌ها آنتروپی به‌صورت تحلیلی صفر است. در خط دوم، فرمول اصلی $-p \log_2 p - (1-p) \log_2(1-p)$ پیاده‌سازی می‌شود. چون این خط فقط در صورتی اجرا می‌شود که $p \in (0, 1)$، فراخوانی \lr{\texttt{math.log2}} همیشه با ورودی مثبت انجام می‌شود و خطای \lr{math domain error} رخ نمی‌دهد. این پیاده‌سازی، الگوی \lr{defensive programming} را به نمایش می‌گذارد.

\subsection{تحلیل تابع analytical\_skr و ساختار پارامترها}

تابع \lr{\texttt{analytical\_skr}} اصلی‌ترین تابع تحلیلی این ماژول است که نرخ کلید امن را بر اساس پنج پارامتر پروتکل محاسبه می‌کند. این تابع نه پارامتر ورودی دارد: \lr{$q_x$} (\lr{sifting factor})، \lr{$p_s$} (احتمال سیگنال)، \lr{$p_d$} (احتمال دکوی)، \lr{$\mu$} (شدت سیگنال)، \lr{$\nu$} (شدت دکوی)، \lr{$\eta$} (ضریب انتقال کلی)، \lr{$Y_0$} (نرخ دارک کانت)، \lr{$e_d$} (خطای ذاتی) و \lr{$f_{\text{EC}}$} (بازدهی \lr{error correction}). خروجی تابع، یک عدد اعشاری است که نرخ کلید امن را نشان می‌دهد.

\begin{LTR}
	\begin{lstlisting}[language=Python]
		def analytical_skr(q_x, p_s, p_d, mu, nu, eta, Y0, ed, f_ec):
		"""تابع هزینه تحلیلی برای بررسی پارامترها (اصلاح شده)"""
		if p_s + p_d >= 1.0: return 0.0
		if nu >= mu * 0.9: return 0.0
		
		Q_mu = 1 - (1 - Y0) * math.exp(-eta * mu)
		E_mu = (0.5 * Y0 + ed * (1 - math.exp(-eta * mu))) / Q_mu
		Q_nu = 1 - (1 - Y0) * math.exp(-eta * nu)
	\end{lstlisting}
\end{LTR}

در خطوط ابتدایی، دو قید فیزیکی بررسی می‌شوند: قید \lr{$p_s + p_d \geq 1$} و قید \lr{$\nu \geq 0.9 \mu$}. در صورت نقض هر یک، تابع بلافاصله مقدار صفر را برمی‌گرداند. این الگوی \lr{early return} از پردازش غیرضروری جلوگیری می‌کند و خوانایی کد را بالا می‌برد. سپس سه کمیت کلیدی $Q_\mu$، $E_\mu$ و $Q_\nu$ محاسبه می‌شوند. این کمیت‌ها پایه محاسبات بعدی هستند.

\begin{LTR}
	\begin{lstlisting}[language=Python]
		if abs(mu - nu) < 1e-9: return 0.0
		
		# اصلاح فرمول Lo2005 برای Y1
		term1 = Q_nu * math.exp(nu)
		term2 = (nu**2 / mu**2) * Q_mu * math.exp(mu)
		factor = mu / (mu * nu - nu**2)
		Y1_L = factor * (term1 - term2)
		
		Y1 = max(Y1_L, Y0)  # Y1 نمی‌تواند کمتر از نویز تاریک باشد
		e1 = (0.5 * Y0 + ed * eta) / Y1 if Y1 > 0 else 0.5
	\end{lstlisting}
\end{LTR}

در این بخش، کران پایین $Y_1^L$ محاسبه می‌شود. شرط $\left| \mu - \nu \right| < 10^{-9}$ حالت تکین را شناسایی می‌کند؛ در این حالت، تابع مقدار صفر را برمی‌گرداند. در غیر این صورت، فرمول اصلی \lr{Lo2005} پیاده می‌شود: \lr{\texttt{term1}} و \lr{\texttt{term2}} دو جمله اصلی صورت کسر هستند و \lr{\texttt{factor}} ضریب نرمال‌سازی است. در نهایت، \lr{\texttt{max(Y1\_L, Y0)}} تضمین می‌کند که $Y_1$ نهایی کم‌تر از $Y_0$ نشود؛ این برش فیزیکی، بازتاب‌دهنده این واقعیت است که حتی در غیاب فوتون واقعی، دارک کانت می‌تواند منجر به کلیک شود.

محاسبه $e_1$ نیز با شرط $Y_1 > 0$ همراه است تا از تقسیم بر صفر جلوگیری شود؛ در صورت $Y_1 = 0$، مقدار $e_1 = 0.5$ گرفته می‌شود که نشان‌دهنده حالت کاملاً تصادفی است.

\begin{LTR}
	\begin{lstlisting}[language=Python]
		Q1 = mu * math.exp(-mu) * Y1
		R = q_x * (Q1 * (1 - H2(e1)) - Q_mu * f_ec * H2(E_mu))
		return max(R, 0.0)
	\end{lstlisting}
\end{LTR}

در بخش پایانی، $Q_1$ و $R$ محاسبه می‌شوند. $Q_1$ از ضرب $\mu e^{-\mu}$ (احتمال پالس تک‌فوتونی) در $Y_1$ (بازده تک‌فوتونی) به دست می‌آید. $R$ از فرمول \lr{GLLP} محاسبه می‌شود و در نهایت، \lr{\texttt{max(R, 0.0)}} تضمین می‌کند که نرخ کلید منفی نشود؛ این برش در رژیم‌های با \lr{QBER} بالا (مثلاً $E_\mu > 0.11$ که آستانه بحرانی \lr{BB84} است) ضروری است، چون در این رژیم، عبارت $Q_\mu \cdot f_{\text{EC}} \cdot H_2(E_\mu)$ بزرگ‌تر از $Q_1 \cdot (1 - H_2(e_1))$ می‌شود و $R$ منفی می‌گردد. در عمل، نرخ کلید منفی به این معناست که هیچ کلید امنی قابل استخراج نیست و مقدار صفر گزارش می‌شود.

\subsection{تحلیل تابع run\_lp\_test و الگوی Cost Callback}

تابع \lr{\texttt{run\_lp\_test}} اصلی‌ترین تابع آزمون است که در یک فاصله مشخص، حلّال بهینه‌سازی را فراخوانی کرده و نتایج را مقایسه می‌نماید. این تابع یک پارامتر ورودی \lr{\texttt{dist\_km}} دارد که فاصله به کیلومتر است.

\begin{LTR}
	\begin{lstlisting}[language=Python]
		def run_lp_test(dist_km):
		print(f"\n[TEST] LP Solver at L={dist_km} km")
		eta_det = 0.15
		Y0 = 3e-6
		ed = 0.01
		f_ec = 1.16
		alpha = 0.2
		eta = 10 ** (-alpha * dist_km / 10) * eta_det
	\end{lstlisting}
\end{LTR}

در ابتدا، پارامترهای فیزیکی ثابت تعریف می‌شوند: $\eta_{\text{det}} = 0.15$، $Y_0 = 3 \times 10^{-6}$، $e_d = 0.01$، $f_{\text{EC}} = 1.16$ و $\alpha = 0.2$ (\lr{dB/km}). سپس ضریب انتقال کلی $\eta$ بر اساس قانون \lr{Beer-Lambert} محاسبه می‌شود. این پارامترها با مقادیر مرجع در ادبیات \lr{QKD} هم‌خوان هستند و در شرایط عملیاتی واقعی برای \lr{QKD} در فواصل مخابراتی متوسط قابل دستیابی هستند.

\begin{LTR}
	\begin{lstlisting}[language=Python]
		params = load_lim2014_dedicated_params(distance_km=dist_km, num_bits=10**7)
		
		def cost_callback(q_x, p_s, p_d, mu, nu):
		# بهینه‌ساز به دنبال مینیمم کردن هزینه است، پس منفی SKR را برمی‌گردانیم
		skr = analytical_skr(q_x, p_s, p_d, mu, nu, eta, Y0, ed, f_ec)
		return -skr
	\end{lstlisting}
\end{LTR}

در این بخش، پارامترهای پروتکل \lr{Lim-2014} برای فاصله و تعداد بیت مشخص بارگذاری می‌شوند. سپس تابع \lr{\texttt{cost\_callback}} به‌عنوان یک \lr{closure} تعریف می‌شود. این تابع، پنج پارامتر پروتکل را می‌گیرد و \emph{منفی} \lr{SKR} را برمی‌گرداند. دلیل این negate کردن آن است که حلّال بهینه‌سازی به‌طور پیش‌فرض به دنبال مینیمم‌کردن تابع هزینه است؛ بنابراین برای حداکثرسازی \lr{SKR}، باید منفی آن را مینیمم کرد. این الگو در ادبیات بهینه‌سازی به‌عنوان \lr{minimization of negated objective} شناخته می‌شود.

این تابع \lr{closure} به متغیرهای \lr{$\eta$}، \lr{$Y_0$}، \lr{$e_d$} و \lr{$f_{\text{EC}}$} از محیط بیرونی دسترسی دارد (\lr{closure capture})؛ این الگو امکان تزریق پارامترهای محیطی به تابع هدف را بدون نیاز به ارسال آن‌ها به‌عنوان پارامتر فراهم می‌سازد. این الگو در توابع \lr{callback} بسیار رایج است و خوانایی کد را بالا می‌برد.

\subsection{الگوی Comparison و Hardcoded Reference}

در ادامه، \lr{SKR} با پارامترهای سخت‌کدشده محاسبه می‌شود و سپس حلّال فراخوانی می‌گردد:

\begin{LTR}
	\begin{lstlisting}[language=Python]
		# ۱. محاسبه SKR با پارامترهای ثابت و سخت‌کد شده
		R_hardcoded = analytical_skr(0.5, 0.5, 0.25, 0.5, 0.1, eta, Y0, ed, f_ec)
		
		# ۲. اجرای LP Solver فریم‌ورک
		try:
		q_x_opt, p_s_opt, p_d_opt, mu_opt, nu_opt = optimize_lim2014_finite_key(
		params, dist_km, cost_callback)
		R_optimized = analytical_skr(q_x_opt, p_s_opt, p_d_opt, mu_opt, nu_opt, 
		eta, Y0, ed, f_ec)
	\end{lstlisting}
\end{LTR}

در این بخش، \lr{$R_{\text{hardcoded}}$} با پارامترهای مرجع \lr{$q_x = 0.5$}، \lr{$p_s = 0.5$}، \lr{$p_d = 0.25$}، \lr{$\mu = 0.5$} و \lr{$\nu = 0.1$} محاسبه می‌شود. این پارامترها، مقادیر کلاسیک و شناخته‌شده در ادبیات \lr{decoy-state QKD} هستند و به‌عنوان نقطه مرجع عمل می‌کنند. سپس حلّال \lr{\texttt{optimize\_lim2014\_finite\_key}} فراخوانی می‌شود و پارامترهای بهینه را برمی‌گرداند. \lr{$R_{\text{optimized}}$} با این پارامترهای بهینه محاسبه می‌شود.

استفاده از بلاک \lr{\texttt{try/except}} برای مدیریت خطاهای احتمالی حلّال ضروری است؛ چون حلّال بهینه‌سازی ممکن است در شرایط دشوار (مثلاً در فواصل بسیار بزرگ که \lr{SKR} نزدیک به صفر است) به همگرایی نرسد یا خطا تولید کند. در صورت وقوع خطا، بلاک \lr{\texttt{except}} پیام خطا و \lr{traceback} کامل را چاپ می‌نماید و تابع مقدار \lr{False} برمی‌گرداند.

\subsection{الگوریتم Comparison و Verdict}

در نهایت، نتایج مقایسه شده و \lr{verdict} اعلام می‌گردد:

\begin{LTR}
	\begin{lstlisting}[language=Python]
		print(f"  Hardcoded (mu=0.5, nu=0.1)     -> SKR: {R_hardcoded:.6e}")
		print(f"  Optimized  (mu={mu_opt:.3f}, nu={nu_opt:.3f}) -> SKR: {R_optimized:.6e}")
		
		if R_optimized >= R_hardcoded - 1e-9:
		print("  PASS: LP Solver successfully found optimal or equal parameters.")
		return True
		else:
		print("  FAIL: LP Solver performed worse than hardcoded parameters!")
		return False
	\end{lstlisting}
\end{LTR}

در این الگوریتم، شرط \lr{$R_{\text{optimized}} \geq R_{\text{hardcoded}} - 10^{-9}$} بررسی می‌شود. در صورت برقراری، آزمون موفق تلقی می‌گردد؛ در غیر این صورت، شکست. استفاده از \lr{$10^{-9}$} به‌عنوان تلورانس، به این دلیل است که در محاسبات عددی، خطای گردکردن می‌تواند منجر به تفاوت‌های بسیار کوچک شود و این تلورانس تضمین می‌کند که چنین تفاوت‌های جزئی به‌عنوان شکست تلقی نشوند. این الگو، الگوی \lr{tolerance-based comparison} است که در آزمون‌های عددی استاندارد است.

\subsection{نقطه ورود \_\_main\_\_ و سه‌آزمون}

\begin{LTR}
	\begin{lstlisting}[language=Python]
		if __name__ == "__main__":
		print("=" * 60)
		print("LP SOLVER (PARAMETER OPTIMIZATION) VALIDATION")
		print("=" * 60)
		
		# تست در فواصل کوتاه، متوسط و طولانی
		res1 = run_lp_test(0.0)
		res2 = run_lp_test(50.0)
		res3 = run_lp_test(100.0)
		
		print("\n" + "=" * 60)
		if res1 and res2 and res3:
		print("100% LP SOLVER VALIDATION PASSED!")
		else:
		print("LP SOLVER VALIDATION FAILED.")
		print("=" * 60)
	\end{lstlisting}
\end{LTR}

نقطه ورود \lr{\texttt{\_\_main\_\_}} یک سرآیند تزئینی چاپ می‌کند و سپس سه آزمون در فواصل $0$، $50$ و $100$ کیلومتر را اجرا می‌نماید. این سه فاصله، سه رژیم فیزیکی متفاوت را پوشش می‌دهند: فاصله $0$ کیلومتر (بدون کاهیدگی)، فاصله $50$ کیلومتر (کاهیدگی متوسط با $\eta_{\text{channel}} \approx 0.1$) و فاصله $100$ کیلومتر (کاهیدگی شدید با $\eta_{\text{channel}} \approx 0.01$). در صورت موفقیت همه سه آزمون، پیام \lr{100\% LP SOLVER VALIDATION PASSED!} چاپ می‌شود و در غیر این صورت، پیام شکست. این الگو، تضمین می‌کند که حلّال در طیف وسیعی از شرایط عملیاتی به‌طور قابل‌اتکا عمل می‌کند.

\section{پیاده‌سازی ریاضی و الگوریتمی}

\subsection{پیاده‌سازی آنتروپی باینری با برش عددی}

پیاده‌سازی تابع $H_2$ در این ماژول، نمونه‌ای از \lr{defensive programming} در محاسبات عددی است. فرمول تحلیلی $-p \log_2 p - (1-p) \log_2(1-p)$ در نقاط $p = 0$ و $p = 1$ به‌صورت تحلیلی به $0$ میل می‌کند، اما در محاسبات عددی، فراخوانی \lr{\texttt{math.log2(0)}} منجر به \lr{ValueError} می‌شود. برای حل این مشکل، شرط $p \leq 0$ یا $p \geq 1$ پیش از فراخوانی \lr{\texttt{math.log2}} بررسی می‌شود و در این موارد، مقدار $0$ بازگردانده می‌شود. این برش عددی با حد گسسته \lr{L'Hôpital} هم‌خوان است و از نظر ریاضی صحیح است.

\subsection{پیاده‌سازی فرمول Lo2005 با اصلاحات فیزیکی}

پیاده‌سازی فرمول $Y_1^L$ در تابع \lr{\texttt{analytical\_skr}} چندین اصلاح مهم دارد نسبت به فرمول کلاسیک \lr{Lo2005}. نخست، فرمول به فرم زیر بازآرایی شده است:

\begin{equation}
	Y_1^{L} = \frac{\mu}{\mu \nu - \nu^2} \left( Q_\nu \cdot e^{\nu} - \frac{\nu^2}{\mu^2} \cdot Q_\mu \cdot e^{\mu} \right)
\end{equation}

این بازآرایی، معادل فرمول کلاسیک \lr{Lo2005} است، اما با ساختار متفاوت. در فرمول کلاسیک، ضریب نرمال‌سازی به فرم $\frac{\mu^2 e^{-\mu}}{\mu \nu - \mu^2}$ است؛ اما در این پیاده‌سازی، ضریب به فرم $\frac{\mu}{\mu \nu - \nu^2}$ بازنویسی شده است. این بازنویسی، با استفاده از رابطه $\mu^2 - \mu \nu = \mu(\mu - \nu) = -\nu(\mu - \nu) + (\mu^2 - \nu^2)$ و simplification جبری، قابل اثبات است.

دوم، اصلاح فیزیکی \lr{\texttt{max(Y1\_L, Y0)}} اعمال شده است. این اصلاح، تضمین می‌کند که $Y_1$ هرگز کم‌تر از $Y_0$ نشود. این برش، بازتاب‌دهنده این واقعیت فیزیکی است که حتی در غیاب فوتون واقعی، دارک کانت می‌تواند منجر به کلیک شود و بنابراین، بازده تک‌فوتونی نمی‌تواند کم‌تر از نرخ دارک کانت باشد. این اصلاح، در برخی نسخه‌های فرمول \lr{Lo2005} به‌صورت ضمنی اعمال می‌شود، اما در این پیاده‌سازی به‌صورت صریح و قابل مشاهده اعمال گردیده است.

سوم، در محاسبه $e_1$، در صورت $Y_1 = 0$، مقدار $e_1 = 0.5$ گرفته می‌شود. این انتخاب، بازتاب‌دهنده حالت کاملاً تصادفی است (یعنی در غیاب سیگنال واقعی، نتیجه اندازه‌گیری کاملاً تصادفی و با احتمال برابر برای هر بیت است). این انتخاب، در طراحی امنیتی \lr{QKD} محافظه‌کارانه است؛ چون فرض می‌کند که در صورت عدم اطمینان، حاملهور می‌تواند کل اطلاعات را بداند.

\subsection{پیاده‌سازی قیدهای فیزیکی و Early Return}

پیاده‌سازی قیدهای فیزیکی با الگوی \lr{early return} انجام می‌شود:

\begin{LTR}
	\begin{lstlisting}[language=Python]
		if p_s + p_d >= 1.0: return 0.0
		if nu >= mu * 0.9: return 0.0
		...
		if abs(mu - nu) < 1e-9: return 0.0
	\end{lstlisting}
\end{LTR}

این الگو، چندین مزیت دارد: نخست، خوانایی کد را بالا می‌برد، چون هر قید به‌صورت مستقل و در ابتدای تابع بررسی می‌شود. دوم، از پردازش غیرضروری جلوگیری می‌کند؛ در صورت نقض قید، بقیه محاسبات انجام نمی‌شود. سوم، در توابع بهینه‌سازی حیاتی است، چون حلّال ممکن است در جریان جستجوی فضای پارامتر، به مقادیر غیرفیزیکی برسد و بدون این قیدها، تابع هدف می‌تواند مقادیر غیرمتعارف یا \lr{NaN} تولید کند که حلّال را در همگرایی مختل می‌نماید.

قید \lr{$\nu \geq 0.9 \mu$} به‌طور خاص جالب توجه است؛ این قید، تلورانسی برای حالت تکین فرمول \lr{Lo2005} فراهم می‌سازد. در فرمول کلاسیک، تکینگی در $\nu = \mu$ رخ می‌دهد، اما در عمل، حتی نزدیک بودن $\nu$ به $\mu$ (مثلاً $\nu = 0.95 \mu$) می‌تواند منجر به نتایج غیرقابل اعتماد شود. این قید، تضمین می‌کند که $\nu$ به‌طور قابل‌توجهی کم‌تر از $\mu$ باشد.

\subsection{الگوریتم Cost Callback و Negate Trick}

الگوریتم \lr{cost callback} با \lr{negate trick} یک الگوی کلاسیک در بهینه‌سازی است:

\begin{LTR}
	\begin{lstlisting}[language=Python]
		def cost_callback(q_x, p_s, p_d, mu, nu):
		skr = analytical_skr(q_x, p_s, p_d, mu, nu, eta, Y0, ed, f_ec)
		return -skr
	\end{lstlisting}
\end{LTR}

در این الگوریتم، تابع هدف واقعی \lr{SKR} محاسبه می‌شود، اما به‌جای بازگرداندن آن، منفی آن بازگردانده می‌شود. دلیل این کار آن است که اکثر حلّال‌های بهینه‌سازی (مانند \lr{scipy.optimize.minimize}) به‌طور پیش‌فرض به دنبال مینیمم‌کردن تابع هستند؛ بنابراین برای حداکثرسازی \lr{SKR}، باید منفی آن را مینیمم کرد. این الگو در ادبیات بهینه‌سازی به‌عنوان \lr{minimization of negated objective} شناخته می‌شود.

این الگو، یک ویژگی مهم دارد: امکان استفاده از همان حلّال برای هر دو مسئله مینیمم‌سازی و حداکثرسازی. به‌عنوان مثال، اگر بخواهیم \lr{QBER} را مینیمم کنیم، تابع هدف \lr{QBER} را مستقیماً برمی‌گردانیم؛ و اگر بخواهیم \lr{SKR} را حداکثر کنیم، منفی آن را برمی‌گردانیم. این یکپارچگی، در طراحی ماژولار سیستم بهینه‌سازی ارزشمند است.

\subsection{الگوریتم Closure Capture برای پارامترهای محیطی}

الگوریتم \lr{closure capture} در تعریف \lr{\texttt{cost\_callback}} به‌کار رفته است:

\begin{LTR}
	\begin{lstlisting}[language=Python]
		eta = 10 ** (-alpha * dist_km / 10) * eta_det
		...
		def cost_callback(q_x, p_s, p_d, mu, nu):
		skr = analytical_skr(q_x, p_s, p_d, mu, nu, eta, Y0, ed, f_ec)
		return -skr
	\end{lstlisting}
\end{LTR}

در این الگوریتم، متغیرهای \lr{$\eta$}، \lr{$Y_0$}، \lr{$e_d$} و \lr{$f_{\text{EC}}$} در محیط بیرونی تابع \lr{\texttt{cost\_callback}} تعریف می‌شوند و تابع \lr{closure} به آن‌ها دسترسی دارد. این الگو، چندین مزیت دارد: نخست، امضای تابع \lr{callback} ساده می‌ماند (فقط پنج پارامتر پروتکل) و حلّال نیازی به دانستن پارامترهای محیطی ندارد. دوم، پارامترهای محیطی می‌توانند در زمان تعریف \lr{closure} محاسبه شوند و در تمام فراخوانی‌های بعدی استفاده گردند؛ این کار از محاسبه مکرر آن‌ها جلوگیری می‌کند. سوم، در صورت نیاز به تغییر پارامترهای محیطی (مثلاً برای آزمون در فاصله دیگر)، فقط کافی است \lr{closure} در یک محیط جدید تعریف شود.

این الگو در برنامه‌نویسی تابعی به‌کررات به‌کار می‌رود و در پایتون با کلمات کلیدی \lr{\texttt{lambda}} یا \lr{\texttt{def}} در داخل تابع دیگر، قابل پیاده‌سازی است.

\subsection{الگوریتم Hardcoded Reference Comparison}

الگوریتم مقایسه با مرجع سخت‌کدشده، الگوریتم ساده اما مؤثری است:

\begin{LTR}
	\begin{lstlisting}[language=Python]
		R_hardcoded = analytical_skr(0.5, 0.5, 0.25, 0.5, 0.1, eta, Y0, ed, f_ec)
		...
		R_optimized = analytical_skr(q_x_opt, p_s_opt, p_d_opt, mu_opt, nu_opt, 
		eta, Y0, ed, f_ec)
		...
		if R_optimized >= R_hardcoded - 1e-9:
		return True
		else:
		return False
	\end{lstlisting}
\end{LTR}

در این الگوریتم، \lr{$R_{\text{hardcoded}}$} با پارامترهای مرجع محاسبه می‌شود و سپس با \lr{$R_{\text{optimized}}$} مقایسه می‌گردد. شرط موفقیت، آن است که \lr{$R_{\text{optimized}}$} حداقل برابر با \lr{$R_{\text{hardcoded}}$} (با تلورانس \lr{$10^{-9}$}) باشد. این الگوریتم، چندین ویژگی دارد:

نخست، استفاده از پارامترهای مرجع \lr{$\mu = 0.5$} و \lr{$\nu = 0.1$} که در ادبیات \lr{decoy-state QKD} به‌عنوان مقادیر کلاسیک شناخته شده‌اند. این پارامترها، نقطه عملیاتی متعارف هستند و در بسیاری از پیاده‌سازی‌های آزمایشگاهی به‌کار رفته‌اند. انتخاب آن‌ها به‌عنوان مرجع، تضمین می‌کند که حلّال حداقل به اندازه یک راه‌حل شناخته‌شده عمل می‌کند.

دوم، استفاده از تلورانس \lr{$10^{-9}$} برای جلوگیری از خطای گردکردن. در محاسبات عددی، خطای گردکردن می‌تواند منجر به تفاوت‌های بسیار کوچک شود و این تلورانس تضمین می‌کند که چنین تفاوت‌های جزئی به‌عنوان شکست تلقی نشوند.

سوم، این الگوریتم، الگوی \lr{necessary condition testing} است: موفقیت آزمون تضمین می‌کند که حلّال حداقل به اندازه مرجع عمل می‌کند؛ هرچند ممکن است پارامترهای بهتر نیز یافت کند. در صورت شکست، قطعاً حلّال دارای باگ است و نیاز به بررسی دارد.

\subsection{الگوریتم Three-Range Testing}

الگوریتم سه‌آزمون در فواصل $0$، $50$ و $100$ کیلومتر، الگوریتم هوشمندانه‌ای است:

\begin{LTR}
	\begin{lstlisting}[language=Python]
		res1 = run_lp_test(0.0)
		res2 = run_lp_test(50.0)
		res3 = run_lp_test(100.0)
	\end{lstlisting}
\end{LTR}

این سه فاصله، سه رژیم فیزیکی متفاوت را پوشش می‌دهند:

نخست، فاصله $L = 0$ کیلومتر: در این حالت، $\eta_{\text{channel}} = 1$ و $\eta_{\text{total}} = \eta_{\text{det}} = 0.15$. در این رژیم، کاهیدگی فیبر وجود ندارد و حلّال باید پارامترهایی با \lr{SKR} نزدیک به حد نظری پیدا کند. در این رژیم، مقدار بهینه $\mu$ معمولاً بزرگ‌تر از $0.5$ است، چون کاهیدگی کم است و پالس‌های چندفوتونی کم‌تر مشکل‌ساز می‌شوند.

دوم، فاصله $L = 50$ کیلومتر: در این حالت، $\eta_{\text{channel}} \approx 0.1$ و $\eta_{\text{total}} \approx 0.015$. در این رژیم، کاهیدگی متوسط است و حلّال باید با کاهش $\mu$ برای جبران کاهیدگی، نرخ کلید را بهینه نماید. مقدار بهینه $\mu$ معمولاً در بازه $0.3$ تا $0.5$ است.

سوم، فاصله $L = 100$ کیلومتر: در این حالت، $\eta_{\text{channel}} \approx 0.01$ و $\eta_{\text{total}} \approx 0.0015$. در این رژیم، کاهیدگی شدید است و حلّال در رژیم نزدیک به آستانه بحرانی \lr{BB84} (\lr{$Q \approx 0.11$}) عمل می‌کند. مقدار بهینه $\mu$ معمولاً در بازه $0.1$ تا $0.3$ است و انتخاب پارامترها به‌طور حساس به نتیجه تأثیر می‌گذارد.

این سه آزمون، تضمین می‌کنند که حلّال در طیف وسیعی از شرایط عملیاتی به‌طور قابل‌اتکا عمل می‌کند و در رژیم‌های مختلف فیزیکی، پارامترهای بهینه را پیدا می‌نماید.

\section{ارسال و دریافت با سایر ماژول‌ها}

\subsection{وابستگی به qkd.proofs.optimization}

اصلی‌ترین وابستگی این ماژول، به \lr{\texttt{qkd.proofs.optimization}} است. این وابستگی از طریق تابع \lr{\texttt{optimize\_lim2014\_finite\_key}} برقرار می‌شود که حلّال بهینه‌سازی پارامترهای پروتکل \lr{Lim-2014} است. این تابع، سه پارامتر ورودی دارد: \lr{\texttt{params}} (پارامترهای پروتکل)، \lr{\texttt{dist\_km}} (فاصله) و \lr{\texttt{cost\_callback}} (تابع هزینه). خروجی آن، پنج پارامتر بهینه (\lr{$q_x$}، \lr{$p_s$}، \lr{$p_d$}، \lr{$\mu$}، \lr{$\nu$}) است.

این وابستگی، یک وابستگی \lr{loose coupling} از طریق \lr{callback} است: حلّال به‌جز امضای تابع \lr{callback}، هیچ چیز دیگری از تابع هدف نمی‌داند. این الگو، امکان جایگزینی تابع هدف با توابع دیگر را فراهم می‌سازد. به‌عنوان مثال، می‌توان به‌جای تابع \lr{analytical\_skr} از یک تابع هزینه پیچیده‌تر (مثلاً \lr{SKR} متناهی-\lr{key} یا تابع هزینه با جریمه‌های اضافی) استفاده کرد، بدون اینکه حلّال نیاز به تغییر داشته باشد.

\subsection{وابستگی به qkd.params}

دومین وابستگی این ماژول، به \lr{\texttt{qkd.params}} است. این وابستگی از طریق تابع \lr{\texttt{load\_lim2014\_dedicated\_params}} برقرار می‌شود که پارامترهای اختصاصی پروتکل \lr{Lim-2014} را بارگذاری می‌نماید. این تابع، دو پارامتر ورودی دارد: \lr{\texttt{distance\_km}} و \lr{\texttt{num\_bits}}. خروجی آن، یک شیء \lr{dict} یا \lr{dataclass} است که شامل پارامترهای امنیتی پروتکل (مانند \lr{$\varepsilon_{\text{PE}}$}، \lr{$\varepsilon_{\text{EC}}$}، \lr{$\varepsilon_{\text{s}}$} و سایر پارامترهای متناهی-\lr{key}) است.

این پارامترها، در حلّال بهینه‌سازی برای محاسبه جریمه‌های متناهی-\lr{key} استفاده می‌شوند. در این ماژول اعتبارسنجی، از این پارامترها فقط برای فراخوانی حلّال استفاده می‌شود و در تابع هدف تحلیلی (\lr{\texttt{analytical\_skr}}) مستقیماً به‌کار نمی‌روند؛ چون تابع هدف، حد مجانبی را محاسبه می‌کند و جریمه‌های متناهی-\lr{key} را در بر نمی‌گیرد. این جداسازی، الگوی \lr{independent oracle} است که در بخش بعد توضیح داده می‌شود.

\subsection{الگوی Independent Oracle و جلوگیری از Circular Validation}

یکی از ویژگی‌های برجسته این ماژول، استفاده از یک تابع هدف تحلیلی مستقل برای اعتبارسنجی است. تابع \lr{\texttt{analytical\_skr}} در این ماژول، مستقل از حلّال پیاده‌سازی شده و به‌عنوان یک \lr{oracle} عمل می‌کند که صحت خروجی حلّال را راستی‌آزمایی می‌نماید. این جداسازی، از خطای \lr{circular validation} جلوگیری می‌کند.

\lr{circular validation} زمانی رخ می‌دهد که اعتبارسنجی از همان تابعی استفاده کند که در پیاده‌سازی اصلی به‌کار رفته است. در این حالت، اگر تابع دارای باگ باشد، این باگ در اعتبارسنجی نیز منتقل می‌شود و شناسایی نمی‌گردد. به‌عنوان مثال، اگر حلّال از یک تابع هدف باگ‌دار (مثلاً با فرمول اشتباه برای $Y_1$) استفاده کند و اعتبارسنجی نیز از همان تابع استفاده نماید، هر دو خروجی یکسانی تولید می‌کنند و آزمون موفق تلقی می‌شود، در حالی که خروجی هر دو نادرست است.

برای جلوگیری از این خطا، تابع \lr{\texttt{analytical\_skr}} در این ماژول، مستقل از حلّال پیاده‌سازی شده است. این تابع، فرمول \lr{GLLP} را به‌صورت تحلیلی و با اصلاحات فیزیکی (مانند \lr{\texttt{max(Y1\_L, Y0)}}) پیاده می‌سازد و به‌عنوان یک مرجع مستقل عمل می‌نماید. در صورت بروز باگ در حلّال، این تابع مستقل، خطا را آشکار می‌سازد.

\subsection{خروجی به stdout و گزارش متمرکز}

خروجی این ماژول به‌صورت چاپ به \lr{stdout} است. این خروجی شامل چندین بخش است: سرآیند تزئینی، نتیجه هر آزمون (شامل پارامترهای سخت‌کدشده و بهینه و \lr{SKR} مربوطه)، و \lr{verdict} نهایی. این قالب خروجی، برای مصرف توسط توسعه‌دهنده در حین \lr{debug} طراحی شده و امکان تحلیل سریع نتایج را فراهم می‌سازد.

در صورت نیاز، می‌توان خروجی را به فایل نیز نوشت تا در \lr{CI/CD pipeline} قابل بازبینی باشد. در نسخه فعلی، این قابلیت پیاده‌سازی نشده است، اما با افزودن چند خط کد (مانند \lr{\texttt{with open(...) as f: f.write(...)}}) قابل افزودن است.

\subsection{ارتباط با absolute\_crypto\_validation.py}

این ماژول و \lr{\texttt{absolute\_crypto\_validation.py}} دو لایه مکمل اعتبارسنجی را تشکیل می‌دهند. \lr{\texttt{absolute\_crypto\_validation.py}} بر راستی‌آزمایی کمیت‌های پایه (\lr{$Q_\mu$}، \lr{$E_\mu$}، \lr{$Y_1$}، \lr{$e_1$}، \lr{$R$}) در برابر فرمول‌های تحلیلی متمرکز است، در حالی که این ماژول (\lr{\texttt{lp\_validation.py}}) بر راستی‌آزمایی حلّال بهینه‌سازی پارامترها متمرکز است. این دو لایه، نقش‌های مکمل دارند: لایه اول تضمین می‌کند که موتور محاسباتی به‌درستی پیاده شده، و لایه دوم تضمین می‌نماید که حلّال بهینه‌سازی به‌درستی عمل می‌کند. در صورت شکست هر یک از این لایه‌ها، مشکلات متفاوتی را آشکار می‌سازند: شکست لایه اول نشان‌دهنده باگ در محاسبات پایه است، در حالی که شکست لایه دوم نشان‌دهنده باگ در حلّال بهینه‌سازی (مثلاً عدم همگرایی یا یافتن کمینه محلی) است.

\subsection{ارتباط با ماژول‌های اثبات Lim-2014 و Ma-2005}

این ماژول به‌طور غیرمستقیم با ماژول‌های اثبات امنیتی \lr{\texttt{lim2014.py}} و \lr{\texttt{ma2005.py}} در ارتباط است. از طریق تابع \lr{\texttt{optimize\_lim2014\_finite\_key}}، این ماژول از پارامترهای پروتکل \lr{Lim-2014} استفاده می‌کند. با این حال، تابع هدف تحلیلی (\lr{\texttt{analytical\_skr}}) بر اساس حد مجانبی است و با فرمول‌های متناهی-\lr{key} در \lr{Lim-2014} متفاوت است. این تفاوت، عمدی است: آزمون کران بالای نرخ کلید را راستی‌آزمایی می‌کند و در صورتی که حلّال نتواند به کران بالا نزدیک شود، قطعاً نمی‌تواند نرخ کلید متناهی صحیح را تولید کند. این یک الگوی \lr{necessary condition testing} است.

\subsection{خروجی نهایی و استفاده در CI/CD}

خروجی نهایی این ماژول به دو شکل است: نخست، خروجی چاپی به کنسول که شامل نتایج هر آزمون و \lr{verdict} نهایی است. دوم، کد خروج \lr{POSIX} (به‌طور ضمنی از طریق \lr{\texttt{sys.exit}}) که در صورت شکست، غیرصفر است. در پیاده‌سازی فعلی، کد خروج به‌صورت صریح تنظیم نشده است، اما می‌توان با افزودن \lr{\texttt{sys.exit(0 if all([res1, res2, res3]) else 1)}} در پایان \lr{\texttt{\_\_main\_\_}}، این قابلیت را اضافه کرد.

در یک \lr{CI/CD pipeline} نمونه، این ماژول به‌صورت زیر به‌کار می‌رود:

\begin{LTR}
	\begin{lstlisting}[language=bash]
		- name: LP Solver Validation
		run: python lp_validation.py
	\end{lstlisting}
\end{LTR}

در صورت شکست، \lr{pipeline} متوقف می‌شود و توسعه‌دهنده باید قبل از \lr{merge}، خطا را اصلاح کند. این الگو، تضمین می‌کند که هرگز کدی با خطای بهینه‌سازی به شاخه اصلی \lr{merge} نشود.

\subsection{جمع‌بندی ارتباطات}

به‌طور خلاصه، این ماژول در نقش یک \lr{validator} مستقل عمل می‌کند که از طریق سه کانال با سایر بخش‌های فریم‌ورک در ارتباط است: ورودی حلّال بهینه‌سازی از \lr{\texttt{qkd.proofs.optimization}}، ورودی پارامترهای پروتکل از \lr{\texttt{qkd.params}}، و خروجی به \lr{stdout} برای مصرف توسعه‌دهنده. این معماری، جداسازی واضحی بین لایه اعتبارسنجی و موتور بهینه‌سازی برقرار می‌سازد و امکان توسعه مستقل هر یک را فراهم می‌نماید. خروجی این ماژول، نه‌تنها برای توسعه‌دهندگان (در فرآیند \lr{debug})، بلکه برای بازبینان کریپتوگرافیک (در فرآیند \lr{security audit}) نیز مفید است، زیرا یک \lr{audit trail} قابل ممیزی از عملکرد حلّال بهینه‌سازی در رژیم‌های مختلف فیزیکی فراهم می‌سازد.
